% Actual source scan PDF335--344, printed323--332. Native bibliographic reconstruction.
% Person names and source work titles remain in their original script.
\begin{thebibliography}{EGAIII2}
\addcontentsline{toc}{chapter}{参考文献}
\small
% PDF335 / printed323
\bibitem[AK70]{AK70} Allen Altman and Steven Kleiman, \emph{Introduction to Grothendieck duality theory}, Lecture Notes in Mathematics, Vol. 146, Springer-Verlag, Berlin, 1970.
\bibitem[AK74]{AK74} \rule{2em}{.3pt}, \emph{Algebraic systems of linearly equivalent divisor-like subschemes}, Compositio Math. \textbf{29} (1974), 113--139.
\bibitem[AK79]{AK79} \rule{2em}{.3pt}, \emph{Compactifying the Picard scheme. II}, Amer. J. Math. \textbf{101} (1979), no. 1, 10--41.
\bibitem[AK80]{AK80} \rule{2em}{.3pt}, \emph{Compactifying the Picard scheme}, Adv. in Math. \textbf{35} (1980), no. 1, 50--112.
\bibitem[AM69]{AM69} M. F. Atiyah and I. G. Macdonald, \emph{Introduction to commutative algebra}, Addison-Wesley Publishing Co., Reading, Mass.-London-Don Mills, Ont., 1969.
\bibitem[Art69a]{Art69a} Michael Artin, \emph{Algebraization of formal moduli. I}, Global Analysis (Papers in Honor of K. Kodaira), Univ. Tokyo Press, Tokyo, 1969, 第 21--71 页.
\bibitem[Art69b]{Art69b} \rule{2em}{.3pt}, \emph{The implicit function theorem in algebraic geometry}, Algebraic Geometry (Internat. Colloq., Tata Inst. Fund. Res., Bombay, 1968), Oxford Univ. Press, London, 1969, 第 13--34 页.
\bibitem[Art71]{Art71} \rule{2em}{.3pt}, \emph{Algebraic spaces}, Yale University Press, New Haven, Conn., 1971, 1969 年在 Yale University 讲授的 A James K. Whittemore Lecture in Mathematics, Yale Mathematical Monographs, 3.
\bibitem[Art73]{Art73} \rule{2em}{.3pt}, \emph{Théorèmes de représentabilité pour les espaces algébriques}, Les Presses de l'Université de Montréal, Montreal, Quebec, 1973, 与 Alexandru Lascu 和 Jean-François Boutot 合作, Séminaire de Mathématiques Supérieures, No. 44 (Été, 1970).
\bibitem[Art74a]{Art74a} \rule{2em}{.3pt}, \emph{Lectures on deformations of singularities}, Tata Institute of Fundamental Research, 1974.
\bibitem[Art74b]{Art74b} \rule{2em}{.3pt}, \emph{Versal deformations and algebraic stacks}, Invent. Math. \textbf{27} (1974), 165--189.
\bibitem[Ati57]{Ati57} M. F. Atiyah, \emph{Complex analytic connections in fibre bundles}, Trans. Amer. Math. Soc. \textbf{85} (1957), 181--207.
\bibitem[Bat99]{Bat99} Victor V. Batyrev, \emph{Birational Calabi-Yau $n$-folds have equal Betti numbers}, New trends in algebraic geometry (Warwick, 1996), London Math. Soc. Lecture Note Ser., vol. 264, Cambridge Univ. Press, Cambridge, 1999, 第 1--11 页.
\bibitem[BB73]{BB73} A. Białynicki-Birula, \emph{Some theorems on actions of algebraic groups}, Ann. of Math. (2) \textbf{98} (1973), 480--497.
\bibitem[BBD82]{BBD82} A. A. Beilinson, J. Bernstein, and P. Deligne, \emph{Faisceaux pervers}, Analysis and topology on singular spaces, I (Luminy, 1981), Astérisque, vol. 100, Soc. Math. France, Paris, 1982, 第 5--171 页.
\bibitem[BCP04]{BCP04} Aldo Brigaglia, Ciro Ciliberto, and Claudio Pedrini, \emph{The Italian school of algebraic geometry and Abel's legacy}, The legacy of Niels Henrik Abel, Springer, Berlin, 2004, 第 295--347 页.
\bibitem[Bea83]{Bea83} Arnaud Beauville, \emph{Variétés Kähleriennes dont la première classe de Chern est nulle}, J. Differential Geom. \textbf{18} (1983), no. 4, 755--782 (1984).
% PDF336 / printed324
\bibitem[BLR90]{BLR90} Siegfried Bosch, Werner Lütkebohmert, and Michel Raynaud, \emph{Néron models}, Ergebnisse der Mathematik und ihrer Grenzgebiete (3), vol. 21, Springer-Verlag, Berlin, 1990.
\bibitem[BLR00]{BLR00} Jean-Benoît Bost, François Loeser, and Michel Raynaud (编), \emph{Courbes semi-stables et groupe fondamental en géométrie algébrique}, Progress in Mathematics, vol. 187, Basel, Birkhäuser Verlag, 2000.
\bibitem[Bor69]{Bor69} Armand Borel, \emph{Linear algebraic groups}, Hyman Bass 记录, W. A. Benjamin, Inc., New York-Amsterdam, 1969.
\bibitem[Bor94a]{Bor94a} Francis Borceux, \emph{Handbook of categorical algebra. 1}, Encyclopedia of Mathematics and its Applications, vol. 50, Cambridge University Press, Cambridge, 1994, Basic category theory.
\bibitem[Bor94b]{Bor94b} \rule{2em}{.3pt}, \emph{Handbook of categorical algebra. 2}, Encyclopedia of Mathematics and its Applications, vol. 51, Cambridge University Press, Cambridge, 1994, Categories and structures.
\bibitem[Bor94c]{Bor94c} \rule{2em}{.3pt}, \emph{Handbook of categorical algebra. 3}, Encyclopedia of Mathematics and its Applications, vol. 52, Cambridge University Press, Cambridge, 1994, Categories of sheaves.
\bibitem[Bou61]{Bou61} Nicolas Bourbaki, \emph{Éléments de mathématique. Fascicule XXVIII. Algèbre commutative. Chapitre 3: Graduations, filtra- tions et topologies. Chapitre 4: Idéaux premiers associés et décomposition primaire}, Actualités Scientifiques et Industrielles, No. 1293, Hermann, Paris, 1961.
% source-error: scan prints "filtra- tions" within the title; retained.
\bibitem[Bou64]{Bou64} \rule{2em}{.3pt}, \emph{Éléments de mathématique. Fasc. XXX. Algèbre commutative. Chapitre 5: Entiers. Chapitre 6: Valuations}, Actualités Scientifiques et Industrielles, No. 1308, Hermann, Paris, 1964.
\bibitem[Bou78]{Bou78} Jean-François Boutot, \emph{Schéma de Picard local}, Lecture Notes in Mathematics, vol. 632, Springer, Berlin, 1978.
\bibitem[BR70]{BR70} Jean Bénabou and Jacques Roubaud, \emph{Monades et descente}, C. R. Acad. Sci. Paris Sér. A-B \textbf{270} (1970), A96--A98.
\bibitem[Bri77]{Bri77} Joël Briançon, \emph{Description de $\mathrm{Hilb}^n\mathbb C\{x,y\}$}, Invent. Math. \textbf{41} (1977), no. 1, 45--89.
\bibitem[Car60]{Car60} Pierre Cartier, \emph{Sur un théorème de Snapper}, Bull. Soc. Math. France \textbf{88} (1960), 333--343.
\bibitem[Che96]{Che96} Jan Cheah, \emph{On the cohomology of Hilbert schemes of points}, J. Algebraic Geom. \textbf{5} (1996), no. 3, 479--511.
\bibitem[Cho57]{Cho57} Wei-Liang Chow, \emph{On the projective embedding of homogeneous varieties}, Algebraic geometry and topology. A symposium in honor of S. Lefschetz, Princeton University Press, Princeton, N. J., 1957, 第 122--128 页.
\bibitem[Con02]{Con02} Brian Conrad, \emph{A modern proof of Chevalley's theorem on algebraic groups}, J. Ramanujan Math. Soc. \textbf{17} (2002), no. 1, 1--18.
\bibitem[CS01]{CS01} Pierre Colmez and Jean-Pierre Serre (编), \emph{Correspondance Grothendieck-Serre}, Documents Mathématiques (Paris), Société Mathématique de France, Paris, 2001.
\bibitem[Dan70]{Dan70} V. I. Danilov, \emph{Samuel's conjecture}, Mat. Sb. (N.S.) \textbf{81 (123)} (1970), 132--144.
\bibitem[dCM00]{dCM00} Mark Andrea A. de Cataldo and Luca Migliorini, \emph{The Douady space of a complex surface}, Adv. Math. \textbf{151} (2000), no. 2, 283--312.
\bibitem[dCM02]{dCM02} \rule{2em}{.3pt}, \emph{The Chow groups and the motive of the Hilbert scheme of points on a surface}, J. Algebra \textbf{251} (2002), no. 2, 824--848.
\bibitem[Del68]{Del68} Pierre Deligne, \emph{Théorème de Lefschetz et critères de dégénérescence de suites spectrales}, Inst. Hautes Études Sci. Publ. Math. (1968), no. 35, 259--278.
% PDF337 / printed325
\bibitem[Del73]{Del73} \rule{2em}{.3pt}, \emph{Cohomologie à supports propres}, Théorie des topos et cohomologie étale des schémas. Tome 3, Springer-Verlag, Berlin, 1973, Séminaire de Géométrie Algébrique du Bois-Marie 1963--1964 (SGA 4), M. Artin, A. Grothendieck 和 J. L. Verdier 主持, P. Deligne 和 B. Saint-Donat 参与合作, Lecture Notes in Mathematics, Vol. 305.
\bibitem[Del81]{Del81} \rule{2em}{.3pt}, \emph{Relèvement des surfaces K3 en caractéristique nulle}, Algebraic surfaces (Orsay, 1976--78), Lecture Notes in Math., vol. 868, Springer, Berlin, 1981, Luc Illusie 整理出版, 第 58--79 页.
\bibitem[DG70]{DG70} Michel Demazure and Pierre Gabriel, \emph{Groupes algébriques. Tome I: Géométrie algébrique, généralités, groupes commutatifs}, Masson \& Cie, Éditeur, Paris, 1970, 附 Michiel Hazewinkel 所作 Corps de classes local.
\bibitem[DHVW85]{DHVW85} L. Dixon, J. A. Harvey, C. Vafa, and E. Witten, \emph{Strings on orbifolds}, Nuclear Phys. B \textbf{261} (1985), no. 4, 678--686.
\bibitem[DHVW86]{DHVW86} L. Dixon, J. Harvey, C. Vafa, and E. Witten, \emph{Strings on orbifolds. II}, Nuclear Phys. B \textbf{274} (1986), no. 2, 285--314.
\bibitem[DI87]{DI87} Pierre Deligne and Luc Illusie, \emph{Relèvements modulo $p^2$ et décomposition du complexe de de Rham}, Invent. Math. \textbf{89} (1987), no. 2, 247--270.
\bibitem[DIX68]{DIX68} \emph{Dix exposés sur la cohomologie des schémas}, Advanced Studies in Pure Mathematics, Vol. 3, North-Holland Publishing Co., Amsterdam, 1968.
\bibitem[DL99]{DL99} Jan Denef and François Loeser, \emph{Germs of arcs on singular algebraic varieties and motivic integration}, Invent. Math. \textbf{135} (1999), no. 1, 201--232.
\bibitem[DM69]{DM69} Pierre Deligne and David Mumford, \emph{The irreducibility of the space of curves of given genus}, Inst. Hautes Études Sci. Publ. Math. \textbf{36} (1969), 75--109.
\bibitem[EG00]{EG00} Geir Ellingsrud and Lothar Göttsche, \emph{Hilbert schemes of points and Heisenberg algebras}, School on Algebraic Geometry (Trieste, 1999), ICTP Lect. Notes, vol. 1, Abdus Salam Int. Cent. Theoret. Phys., Trieste, 2000, 第 59--100 页.
\bibitem[EGAI]{EGAI} Alexander Grothendieck, \emph{Éléments de géométrie algébrique. I. Le langage des schémas}, Inst. Hautes Études Sci. Publ. Math. (1960), no. 4, 228.
\bibitem[EGAII]{EGAII} \rule{2em}{.3pt}, \emph{Éléments de géométrie algébrique. II. Étude globale élémentaire de quelques classes de morphismes}, Inst. Hautes Études Sci. Publ. Math. (1961), no. 8, 222.
\bibitem[EGAIII1]{EGAIII1} \rule{2em}{.3pt}, \emph{Éléments de géométrie algébrique. III. Étude cohomologique des faisceaux cohérents. I}, Inst. Hautes Études Sci. Publ. Math. (1961), no. 11, 167.
\bibitem[EGAIII2]{EGAIII2} \rule{2em}{.3pt}, \emph{Éléments de géométrie algébrique. III. Étude cohomologique des faisceaux cohérents. II}, Inst. Hautes Études Sci. Publ. Math. (1963), no. 17, 91.
\bibitem[EGAIV1]{EGAIV1} \rule{2em}{.3pt}, \emph{Éléments de géométrie algébrique. IV. Étude locale des schémas et des morphismes de schémas. I}, Inst. Hautes Études Sci. Publ. Math. (1964), no. 20, 259.
\bibitem[EGAIV2]{EGAIV2} \rule{2em}{.3pt}, \emph{Éléments de géométrie algébrique. IV. Étude locale des schémas et des morphismes de schémas. II}, Inst. Hautes Études Sci. Publ. Math. (1965), no. 24, 231.
\bibitem[EGAIV3]{EGAIV3} \rule{2em}{.3pt}, \emph{Éléments de géométrie algébrique. IV. Étude locale des schémas et des morphismes de schémas. III}, Inst. Hautes Études Sci. Publ. Math. (1966), no. 28, 255.
\bibitem[EGAIV4]{EGAIV4} \rule{2em}{.3pt}, \emph{Éléments de géométrie algébrique. IV. Étude locale des schémas et des morphismes de schémas IV}, Inst. Hautes Études Sci. Publ. Math. (1967), no. 32, 361.
\bibitem[EGAG]{EGAG} \rule{2em}{.3pt}, \emph{Éléments de géométrie algébrique. I}, Grundlehren Math. Wiss., vol. 166, Springer-Verlag, 1971.
\bibitem[EGK02]{EGK02} Eduardo Esteves, Mathieu Gagné, and Steven Kleiman, \emph{Autoduality of the compactified Jacobian}, J. London Math. Soc. (2) \textbf{65} (2002), no. 3, 591--610.
\bibitem[Eke86]{Eke86} Torsten Ekedahl, \emph{Diagonal complexes and $F$-gauge structures}, Travaux en Cours, Hermann, Paris, 1986.
\bibitem[Ell]{Ell} Geir Ellingsrud, \emph{Irreducibility of the punctual Hilbert scheme of a surface}, 未发表.
% PDF338 / printed326
\bibitem[ES87]{ES87} Geir Ellingsrud and Stein Arild Strømme, \emph{On the homology of the Hilbert scheme of points in the plane}, Invent. Math. \textbf{87} (1987), no. 2, 343--352.
\bibitem[ES98]{ES98} \rule{2em}{.3pt}, \emph{An intersection number for the punctual Hilbert scheme of a surface}, Trans. Amer. Math. Soc. \textbf{350} (1998), no. 6, 2547--2552.
\bibitem[Fal03]{Fal03} Gerd Faltings, \emph{Finiteness of coherent cohomology for proper fppf stacks}, J. Algebraic Geom. \textbf{12} (2003), no. 2, 357--366.
\bibitem[FGA]{FGA} Alexander Grothendieck, \emph{Fondements de la géométrie algébrique. Extraits du Séminaire Bourbaki, 1957--1962}, Secrétariat mathématique, Paris, 1962.
\bibitem[FK88]{FK88} Eberhard Freitag and Reinhardt Kiehl, \emph{Étale cohomology and the Weil conjecture}, Ergebnisse der Mathematik und ihrer Grenzgebiete (3), vol. 13, Springer-Verlag, Berlin, 1988.
\bibitem[FK95]{FK95} Kazuhiro Fujiwara and Kazuya Kato, \emph{Logarithmic étale topology theory}, (未完成的) 预印本, 1995.
\bibitem[FM98]{FM98} Barbara Fantechi and Marco Manetti, \emph{Obstruction calculus for functors of Artin rings. I}, J. Algebra \textbf{202} (1998), no. 2, 541--576.
\bibitem[Fog68]{Fog68} John Fogarty, \emph{Algebraic families on an algebraic surface}, Amer. J. Math \textbf{90} (1968), 511--521.
\bibitem[Fos81]{Fos81} Robert Fossum, \emph{Formes différentielles non fermées}, Séminaire sur les Pinceaux de Courbes de Genre au Moins Deux (Lucien Szpiro 编), Astérisque, vol. 86, Société Mathématique de France, Paris, 1981, 第 90--96 页.
\bibitem[Gir64]{Gir64} Jean Giraud, \emph{Méthode de la descente}, Bull. Soc. Math. France Mém. \textbf{2} (1964), viii+150.
\bibitem[Gir71]{Gir71} \rule{2em}{.3pt}, \emph{Cohomologie non abélienne}, Springer-Verlag, Berlin, 1971, Die Grundlehren der mathematischen Wissenschaften, Band 179.
\bibitem[God58]{God58} Roger Godement, \emph{Topologie algébrique et théorie des faisceaux}, Actualit'es Sci. Ind. No. 1252. Publ. Math. Univ. Strasbourg. No. 13, Hermann, Paris, 1958.
% source-error: scan prints Actualit'es, retained.
\bibitem[Göt90]{Göt90} Lothar Göttsche, \emph{The Betti numbers of the Hilbert scheme of points on a smooth projective surface}, Math. Ann. \textbf{286} (1990), no. 1--3, 193--207.
\bibitem[Göt01]{Göt01} \rule{2em}{.3pt}, \emph{On the motive of the Hilbert scheme of points on a surface}, Math. Res. Lett. \textbf{8} (2001), no. 5--6, 613--627.
\bibitem[Göt02]{Göt02} \rule{2em}{.3pt}, \emph{Hilbert schemes of points on surfaces}, Proceedings of the International Congress of Mathematicians, Vol. II (Beijing, 2002) (Beijing), Higher Ed. Press, 2002, 第 483--494 页.
\bibitem[Gra66]{Gra66} John W. Gray, \emph{Fibred and cofibred categories}, Proc. Conf. Categorical Algebra (La Jolla, Calif., 1965), Springer, New York, 1966, 第 21--83 页.
\bibitem[Gro95a]{Gro95a} Alexander Grothendieck, \emph{Géométrie formelle et géométrie algébrique}, Séminaire Bourbaki, Vol. 5, Soc. Math. France, Paris, 1995, Exp. No. 182, 第 193--220 页, 勘误见第 390 页.
\bibitem[Gro95b]{Gro95b} \rule{2em}{.3pt}, \emph{Technique de descente et théorèmes d'existence en géométrie algébrique. I. Généralités. Descente par morphismes fidèlement plats.}, Séminaire Bourbaki, Exp. No. 190, vol. 5, Soc. Math. France, Paris, 1995, 第 299--327 页.
\bibitem[Gro95c]{Gro95c} \rule{2em}{.3pt}, \emph{Techniques de construction et théorèmes d'existence en géométrie algébrique. IV. Les schémas de Hilbert}, Séminaire Bourbaki, Vol. 6, Soc. Math. France, Paris, 1995, Exp. No. 221, 第 249--276 页.
\bibitem[Gro96]{Gro96} I. Grojnowski, \emph{Instantons and affine algebras. I. The Hilbert scheme and vertex operators}, Math. Res. Lett. \textbf{3} (1996), no. 2, 275--291.
\bibitem[GS93]{GS93} Lothar Göttsche and Wolfgang Soergel, \emph{Perverse sheaves and the cohomology of Hilbert schemes of smooth algebraic surfaces}, Math. Ann. \textbf{296} (1993), no. 2, 235--245.
\bibitem[Har66]{Har66} Robin Hartshorne, \emph{Residues and duality}, 1963/64 年在 Harvard 举行的关于 A. Grothendieck 工作的讨论班讲义, 附 P. Deligne 所作附录. Lecture Notes in Mathematics, No. 20, Springer-Verlag, Berlin, 1966.
% PDF339 / printed327
\bibitem[Har70]{Har70} \rule{2em}{.3pt}, \emph{Ample subvarieties of algebraic varieties}, 与 C. Musili 合作撰写的讲义. Lecture Notes in Mathematics, Vol. 156, Springer-Verlag, Berlin, 1970.
\bibitem[Har77]{Har77} \rule{2em}{.3pt}, \emph{Algebraic geometry}, Graduate Texts in Mathematics, no. 52, Springer-Verlag, New York, 1977.
\bibitem[Har92]{Har92} Joe Harris, \emph{Algebraic geometry. a first course}, Graduate Texts in Mathematics, vol. 133, Springer-Verlag, New York, 1992.
\bibitem[Hir62]{Hir62} Heisuke Hironaka, \emph{An example of a non-Kählerian complex-analytic deformation of Kählerian complex structures}, Ann. of Math. (2) \textbf{75} (1962), 190--208.
\bibitem[Hir99]{Hir99} Masayuki Hirokado, \emph{A non-liftable Calabi-Yau threefold in characteristic 3}, Tohoku Math. J. (2) \textbf{51} (1999), no. 4, 479--487.
\bibitem[HS]{HS} André Hirschowitz and Carlos Simpson, \emph{Descente pour les $n$-champs}, arXiv: math.AG/9807049.
\bibitem[Iar72]{Iar72} Anthony Iarrobino, \emph{Reducibility of the families of 0-dimensional schemes on a variety}, Invent. Math. \textbf{15} (1972), 72--77.
\bibitem[Iar77]{Iar77} \rule{2em}{.3pt}, \emph{Punctual Hilbert schemes}, Mem. Amer. Math. Soc. \textbf{10} (1977), no. 188, viii+112.
\bibitem[Iar85]{Iar85} \rule{2em}{.3pt}, \emph{Compressed algebras and components of the punctual Hilbert scheme}, Algebraic geometry, Sitges (Barcelona), 1983, Lecture Notes in Math., vol. 1124, Springer, Berlin, 1985, 第 146--165 页.
\bibitem[Iar87]{Iar87} \rule{2em}{.3pt}, \emph{Hilbert scheme of points: overview of last ten years}, Algebraic geometry, Bowdoin, 1985 (Brunswick, Maine, 1985), Proc. Sympos. Pure Math., vol. 46, Amer. Math. Soc., Providence, RI, 1987, 第 297--320 页.
\bibitem[Igu55]{Igu55} Jun-ichi Igusa, \emph{On some problems in abstract algebraic geometry}, Proc. Nat. Acad. Sci. U. S. A. \textbf{41} (1955), 964--967.
\bibitem[Ill71]{Ill71} Luc Illusie, \emph{Complexe cotangent et déformations. I}, Springer-Verlag, Berlin, 1971, Lecture Notes in Mathematics, Vol. 239.
\bibitem[Ill72a]{Ill72a} \rule{2em}{.3pt}, \emph{Complexe cotangent et déformations. II}, Springer-Verlag, Berlin, 1972, Lecture Notes in Mathematics, Vol. 283.
\bibitem[Ill72b]{Ill72b} \rule{2em}{.3pt}, \emph{Cotangent complex and deformations of torsors and group schemes}, Toposes, algebraic geometry and logic (Conf., Dalhousie Univ., Halifax, N.S., 1971), Springer, Berlin, 1972, 第 159--189 页. Lecture Notes in Math., Vol. 274.
\bibitem[Ill85]{Ill85} \rule{2em}{.3pt}, \emph{Déformations de groupes de Barsotti-Tate (d'après A. Grothendieck)}, Astérisque (1985), no. 127, 151--198, Seminar on arithmetic bundles: the Mordell conjecture (Paris, 1983/84).
\bibitem[Ill02]{Ill02} \rule{2em}{.3pt}, \emph{An overview of the work of K. Fujiwara, K. Kato, and C. Nakayama on logarithmic étale cohomology}, Astérisque (2002), no. 279, 271--322, Cohomologies $p$-adiques et applications arithmétiques, II.
\bibitem[Ive70]{Ive70} Birger Iversen, \emph{Linear determinants with applications to the Picard scheme of a family of algebraic curves}, Springer-Verlag, Berlin, 1970, Lecture Notes in Mathematics, Vol. 174.
\bibitem[Jou79]{Jou79} Jean-Pierre Jouanolou, \emph{Théorèmes de Bertini et applications}, Université Louis Pasteur Département de Mathématique Institut de Recherche Mathématique Avancée, Strasbourg, 1979.
\bibitem[JT84]{JT84} André Joyal and Myles Tierney, \emph{An extension of the Galois theory of Grothendieck}, Mem. Amer. Math. Soc. \textbf{51} (1984), no. 309, vii+71.
\bibitem[Kis00]{Kis00} Mark Kisin, \emph{Prime to $p$ fundamental groups, and tame Galois actions}, Ann. Inst. Fourier (Grenoble) \textbf{50} (2000), no. 4, 1099--1126.
% PDF340 / printed328
\bibitem[Kle71]{Kle71} Steven L. Kleiman, \emph{Les théorèmes de finitude pour le foncteur de Picard}, Théorie des intersections et théorème de Riemann-Roch, Springer-Verlag, 1971, Séminaire de Géométrie Algébrique du Bois-Marie 1966--1967 (SGA 6), P. Berthelot, A. Grothendieck 和 L. Illusie 主持, D. Ferrand, J. P. Jouanolou, O. Jussila, S. Kleiman, M. Raynaud 和 J-P. Serre 参与合作, Lecture Notes in Mathematics, Vol. 225.
\bibitem[Kle73]{Kle73} \rule{2em}{.3pt}, \emph{Completeness of the characteristic system}, Advances in Math. \textbf{11} (1973), 304--310.
\bibitem[Kle04]{Kle04} \rule{2em}{.3pt}, \emph{What is Abel's theorem anyway?}, The legacy of Niels Henrik Abel, Springer, Berlin, 2004, 第 395--440 页.
\bibitem[Knu71]{Knu71} Donald Knutson, \emph{Algebraic spaces}, Springer-Verlag, Berlin, 1971, Lecture Notes in Mathematics, Vol. 203.
\bibitem[KO74]{KO74} Max-Albert Knus and Manuel Ojanguren, \emph{Théorie de la descente et algèbres d'Azumaya}, vol. 389, Springer-Verlag, Berlin, 1974, Lecture Notes in Mathematics.
\bibitem[Kod05]{Kod05} Kunihiko Kodaira, \emph{Complex manifolds and deformation of complex structures}, 英文版, Classics in Mathematics, Springer-Verlag, Berlin, 2005, Kazuo Akao 据 1981 年日文原版译出.
\bibitem[KR04a]{KR04a} Maxim Kontsevich and Alexander Rosenberg, \emph{Noncommutative spaces and flat descent}, Max Planck Institute of Mathematics 预印本, 2004.
\bibitem[KR04b]{KR04b} \rule{2em}{.3pt}, \emph{Noncommutative stacks}, Max Planck Institute of Mathematics 预印本, 2004.
\bibitem[KS04]{KS04} Kazuya Kato and Takeshi Saito, \emph{On the conductor formula of Bloch}, Publ. Math. Inst. Hautes Études Sci. (2004), no. 100, 5--151.
\bibitem[Lan79]{Lan79} William E. Lang, \emph{Quasi-elliptic surfaces in characteristic three}, Ann. Sci. École Norm. Sup. (4) \textbf{12} (1979), no. 4, 473--500.
\bibitem[Leh99]{Leh99} Manfred Lehn, \emph{Chern classes of tautological sheaves on Hilbert schemes of points on surfaces}, Invent. Math. \textbf{136} (1999), no. 1, 157--207.
\bibitem[Leh04]{Leh04} \rule{2em}{.3pt}, \emph{Symplectic moduli spaces}, Intersection Theory and Moduli, ICTP Lect. Notes, Abdus Salam Int. Cent. Theoret. Phys. (Trieste), 2004, 第 139--184 页.
\bibitem[Lip74]{Lip74} Joseph Lipman, \emph{Picard schemes of formal schemes; application to rings with discrete divisor class group}, Classification of algebraic varieties and compact complex manifolds, Springer, Berlin, 1974, 第 94--132 页. Lecture Notes in Math., Vol. 412.
\bibitem[Lip75]{Lip75} \rule{2em}{.3pt}, \emph{Unique factorization in complete local rings}, Algebraic geometry (Proc. Sympos. Pure Math., Vol. 29, Humboldt State Univ., Arcata, Calif., 1974), Amer. Math. Soc., Providence, R.I., 1975, 第 531--546 页.
\bibitem[Lip76]{Lip76} \rule{2em}{.3pt}, \emph{The Picard group of a scheme over an Artin ring}, Inst. Hautes Études Sci. Publ. Math. (1976), no. 46, 15--86.
\bibitem[LLR04]{LLR04} Qing Liu, Dino Lorenzini, and Michel Raynaud, \emph{Néron models, Lie algebras, and reduction of curves of genus one}, Invent. Math. \textbf{157} (2004), no. 3, 455--518.
\bibitem[LMB00]{LMB00} Gérard Laumon and Laurent Moret-Bailly, \emph{Champs algébriques}, Ergebnisse der Mathematik und ihrer Grenzgebiete. 3. Folge., vol. 39, Springer-Verlag, Berlin, 2000.
\bibitem[LN80]{LN80} William E. Lang and Niels O. Nygaard, \emph{A short proof of the Rudakov-Šafarevič theorem}, Math. Ann. \textbf{251} (1980), no. 2, 171--173.
\bibitem[LQW02]{LQW02} Wei-ping Li, Zhenbo Qin, and Weiqiang Wang, \emph{Vertex algebras and the cohomology ring structure of Hilbert schemes of points on surfaces}, Math. Ann. \textbf{324} (2002), no. 1, 105--133.
\bibitem[LS67]{LS67} S. Lichtenbaum and M. Schlessinger, \emph{The cotangent complex of a morphism}, Trans. Amer. Math. Soc. \textbf{128} (1967), 41--70.
\bibitem[LS01]{LS01} Manfred Lehn and Christoph Sorger, \emph{Symmetric groups and the cup product on the cohomology of Hilbert schemes}, Duke Math. J. \textbf{110} (2001), no. 2, 345--357.
% PDF341 / printed329
\bibitem[Mac62]{Mac62} I. G. Macdonald, \emph{The Poincaré polynomial of a symmetric product}, Proc. Cambridge Philos. Soc. \textbf{58} (1962), 563--568.
\bibitem[Mat89]{Mat89} Hideyuki Matsumura, \emph{Commutative ring theory}, 第 2 版, Cambridge Studies in Advanced Mathematics, vol. 8, Cambridge University Press, Cambridge, 1989, M. Reid 由日文译出.
\bibitem[MFK94]{MFK94} David Mumford, John Fogarty, and Frances C. Kirwan, \emph{Geometric invariant theory}, 第 3 版, Ergebnisse der Mathematik und ihrer Grenzgebiete (2), vol. 34, Springer-Verlag, Berlin, 1994.
\bibitem[Mil80]{Mil80} James S. Milne, \emph{Étale cohomology}, Princeton Mathematical Series, vol. 33, Princeton University Press, Princeton, N.J., 1980.
\bibitem[ML98]{ML98} Saunders Mac Lane, \emph{Categories for the working mathematician}, 第 2 版, Graduate Texts in Mathematics, vol. 5, Springer-Verlag, New York, 1998.
\bibitem[Moe02]{Moe02} Ieke Moerdijk, \emph{Introduction to the language of stacks and gerbes}, arXiv: math.AT/0212266, 2002.
\bibitem[Mum]{Mum} David Mumford, \emph{Letter to Serre}, 1961 年 5 月.
\bibitem[Mum61a]{Mum61a} \rule{2em}{.3pt}, \emph{Pathologies of modular algebraic surfaces}, Amer. J. Math. \textbf{83} (1961), 339--342.
\bibitem[Mum61b]{Mum61b} \rule{2em}{.3pt}, \emph{The topology of normal singularities of an algebraic surface and a criterion for simplicity}, Inst. Hautes Études Sci. Publ. Math. (1961), no. 9, 5--22.
\bibitem[Mum65]{Mum65} \rule{2em}{.3pt}, \emph{Geometric invariant theory}, Ergebnisse der Mathematik und ihrer Grenzgebiete, Neue Folge, Band 34, Springer-Verlag, Berlin, 1965.
\bibitem[Mum66]{Mum66} \rule{2em}{.3pt}, \emph{Lectures on curves on an algebraic surface}, 含 G. M. Bergman 所作的一节. Annals of Mathematics Studies, No. 59, Princeton University Press, Princeton, N.J., 1966.
\bibitem[Mum67]{Mum67} \rule{2em}{.3pt}, \emph{Pathologies. III}, Amer. J. Math. \textbf{89} (1967), 94--104.
\bibitem[Mum69]{Mum69} \rule{2em}{.3pt}, \emph{Bi-extensions of formal groups}, Algebraic Geometry (Internat. Colloq., Tata Inst. Fund. Res., Bombay, 1968), Oxford Univ. Press, London, 1969, 第 307--322 页.
\bibitem[Mum70]{Mum70} \rule{2em}{.3pt}, \emph{Abelian varieties}, Tata Institute of Fundamental Research Studies in Mathematics, No. 5, 为 Tata Institute of Fundamental Research, Bombay 出版; Oxford University Press, London, 1970.
\bibitem[Mur64]{Mur64} J. P. Murre, \emph{On contravariant functors from the category of pre-schemes over a field into the category of abelian groups (with an application to the Picard functor)}, Inst. Hautes Études Sci. Publ. Math. (1964), no. 23, 5--43.
\bibitem[Mur95]{Mur95} \rule{2em}{.3pt}, \emph{Representation of unramified functors. Applications (according to unpublished results of A. Grothendieck)}, Séminaire Bourbaki, Vol. 9, Soc. Math. France, Paris, 1995, Exp. No. 294, 第 243--261 页.
\bibitem[Nag56]{Nag56} Masayoshi Nagata, \emph{A general theory of algebraic geometry over Dedekind domains. I. The notion of models}, Amer. J. Math. \textbf{78} (1956), 78--116.
\bibitem[Nak97]{Nak97} Hiraku Nakajima, \emph{Heisenberg algebra and Hilbert schemes of points on projective surfaces}, Ann. of Math. (2) \textbf{145} (1997), no. 2, 379--388.
\bibitem[Naq99]{Naq99} \rule{2em}{.3pt}, \emph{Lectures on Hilbert schemes of points on surfaces}, University Lecture Series, vol. 18, American Mathematical Society, Providence, RI, 1999.
% source-error: the source printed key is Naq99 (q), retained.
\bibitem[Nyg79]{Nyg79} Niels O. Nygaard, \emph{A $p$-adic proof of the nonexistence of vector fields on K3 surfaces}, Ann. of Math. (2) \textbf{110} (1979), no. 3, 515--528.
\bibitem[OB72]{OB72} Arthur Ogus and George Bergman, \emph{Nakayama's lemma for half-exact functors}, Proc. Amer. Math. Soc. \textbf{31} (1972), 67--74.
\bibitem[Ols]{Ols} Martin Olsson, \emph{On proper coverings of artin stacks}, 将发表于 Adv. Math.
\bibitem[Oor62]{Oor62} Frans Oort, \emph{Sur le schéma de Picard}, Bull. Soc. Math. France \textbf{90} (1962), 1--14.
% PDF342 / printed330
\bibitem[Oor71]{Oor71} \rule{2em}{.3pt}, \emph{Finite group schemes, local moduli for abelian varieties, and lifting problems}, 载于 Proc. 5th Nordic Summer School Oslo, 1970, Wolters-Noordhoff, 1972, 第 223--254 页 $\cdot=\cdot$ Compositio Math. \textbf{23} (1971), 265--296.
\bibitem[OV00]{OV00} Fabrice Orgogozo and Isabelle Vidal, \emph{Le théorème de spécialisation du groupe fondamental}, Courbes semi-stables et groupe fondamental en géométrie algébrique (Luminy, 1998), Progr. Math., vol. 187, Birkhäuser, Basel, 2000, 第 169--184 页.
\bibitem[Ray66]{Ray66} Michel Raynaud, \emph{Faisceaux amples sur les schémas en groupes et les espaces homogènes}, C. R. Acad. Sci. Paris Sér. A-B \textbf{262} (1966), A1313--A1315.
\bibitem[Ray70]{Ray70} \rule{2em}{.3pt}, \emph{Faisceaux amples sur les schémas en groupes et les espaces homogènes}, Lecture Notes in Mathematics, no. 119, Springer-Verlag, 1970.
\bibitem[Ray71]{Ray71} \rule{2em}{.3pt}, \emph{Un théorème de représentabilité pour le foncteur de Picard}, Théorie des intersections et théorème de Riemann-Roch, Springer-Verlag, 1971, Exposé XII, S. Kleiman 整理. 载于 Séminaire de Géométrie Algébrique du Bois-Marie 1966--1967 (SGA 6), P. Berthelot, A. Grothendieck 和 L. Illusie 主持, D. Ferrand, J. P. Jouanolou, O. Jussila, S. Kleiman, M. Raynaud 和 J-P. Serre 参与合作, Lecture Notes in Mathematics, Vol. 225.
\bibitem[Ray94]{Ray94} \rule{2em}{.3pt}, \emph{Revêtements de la droite affine en caractéristique $p>0$ et conjecture d'Abhyankar}, Invent. Math. \textbf{116} (1994), no. 1--3, 425--462.
\bibitem[RŠ76]{RŠ76} A. N. Rudakov and I. R. Šafarevič, \emph{Inseparable morphisms of algebraic surfaces}, Izv. Akad. Nauk SSSR Ser. Mat. \textbf{40} (1976), no. 6, 1269--1307, 1439.
\bibitem[Sch68]{Sch68} Michael Schlessinger, \emph{Functors of Artin rings}, Trans. Amer. Math. Soc. \textbf{130} (1968), 208--222.
\bibitem[Seg37]{Seg37} Beniamino Segre, \emph{Intorno ad un teorema di hodge sulla teoria della base per le curve di una superficie algebrica}, Ann. di Mat. pura e applicata \textbf{16} (1937), 157--163.
\bibitem[Sei50]{Sei50} A. Seidenberg, \emph{The hyperplane sections of normal varieties}, Trans. Amer. Math. Soc. \textbf{69} (1950), 357--386.
\bibitem[Ser55]{Ser55} Jean-Pierre Serre, \emph{Faisceaux algébriques cohérents}, Ann. of Math. (2) \textbf{61} (1955), 197--278 ([Ser] 中的第 29 篇).
\bibitem[Ser56]{Ser56} \rule{2em}{.3pt}, \emph{Géométrie algébrique et géométrie analytique}, Ann. Inst. Fourier, Grenoble \textbf{6} (1955--1956), 1--42 ([Ser] 中的第 32 篇).
\bibitem[Ser58]{Ser58} \rule{2em}{.3pt}, \emph{Sur la topologie des variétés algébriques en caractéristique $p$}, Symposium internacional de topología algebraica, Universidad Nacional Autónoma de México and UNESCO, Mexico City, 1958, 第 24--53 页 ([Ser] 中的第 38 篇).
\bibitem[Ser59]{Ser59} \rule{2em}{.3pt}, \emph{Groupes algébriques et corps de classes}, Publications de l'institut de mathématique de l'université de Nancago, VII. Hermann, Paris, 1959.
\bibitem[Ser61]{Ser61} \rule{2em}{.3pt}, \emph{Exemples de variétés projectives en caractéristique $p$ non relevables en caractéristique zéro}, Proc. Nat. Acad. Sci. U.S.A. \textbf{47} (1961), 108--109 ([Ser] 中的第 50 篇).
\bibitem[Ser]{Ser} \rule{2em}{.3pt}, \emph{Oeuvres (Collected Papers)}, I, II, III, Springer-Verlag (1986); IV, Springer-Verlag (2000).
\bibitem[SGA1]{SGA1} Alexander Grothendieck, \emph{Revêtements étales et groupe fondamental (SGA 1)}, Lecture Notes in Math., 224, Springer-Verlag, Berlin, 1964, Séminaire de géométrie algébrique du Bois Marie 1960--61, A. Grothendieck 主持, 含 M. Raynaud 所作的两篇论文.
% source-error: scan prints 1964 for SGA1 here; retained.
\bibitem[SGA2]{SGA2} \rule{2em}{.3pt}, \emph{Cohomologie locale des faisceaux cohérents et théorèmes de Lefschetz locaux et globaux (SGA 2)}, North-Holland Publishing Co., Amsterdam, 1968, 增补 Michèle Raynaud 的一篇报告, Séminaire de Géométrie Algébrique du Bois-Marie, 1962, Advanced Studies in Pure Mathematics, Vol. 2.
\bibitem[SGA4]{SGA4} Michael Artin, Alexander Grothendieck, and Jean-Louis Verdier, \emph{Théorie des topos et cohomologie étale des schémas}, 1, 2, 3, Springer-Verlag, Berlin, 1972--73, Séminaire de Géométrie Algébrique du Bois-Marie 1963--1964 (SGA 4), M. Artin, A. Grothendieck 和 J. L. Verdier 主持, N. Bourbaki, P. Deligne 和 B. Saint-Donat 参与合作, Lecture Notes in Mathematics, Voll. 270, 305, 569.
% source-error: printed SGA4 volume numbers 270,305,569 are retained, not replaced with external data.
% PDF343 / printed331
\bibitem[SGA4-1/2]{SGA4-1/2} Pierre Deligne, \emph{Cohomologie étale}, Springer-Verlag, Berlin, 1977, Séminaire de Géométrie Algébrique du Bois-Marie SGA $4\frac12$, J. F. Boutot, A. Grothendieck, L. Illusie 和 J. L. Verdier 参与合作, Lecture Notes in Mathematics, Vol. 569.
\bibitem[SGA5]{SGA5} \emph{Cohomologie $l$-adique et fonctions $L$}, Springer-Verlag, Berlin, 1977, Séminaire de Géometrie Algébrique du Bois-Marie 1965--1966 (SGA 5), Luc Illusie 编, Lecture Notes in Mathematics, Vol. 589.
\bibitem[SGA6]{SGA6} \emph{Théorie des intersections et théorème de Riemann-Roch}, Springer-Verlag, Berlin, 1971, Séminaire de Géométrie Algébrique du Bois-Marie 1966--1967 (SGA 6), P. Berthelot, A. Grothendieck 和 L. Illusie 主持, D. Ferrand, J. P. Jouanolou, O. Jussila, S. Kleiman, M. Raynaud 和 J-P. Serre 参与合作, Lecture Notes in Mathematics, Vol. 225.
\bibitem[SGA7]{SGA7} \emph{Groupes de monodromie en géométrie algébrique. I}, Springer-Verlag, Berlin, 1972, Séminaire de Géométrie Algébrique du Bois-Marie 1967--1969 (SGA 7 I), A. Grothendieck 主持, M. Raynaud 和 D. S. Rim 参与合作, Lecture Notes in Mathematics, Vol. 288.
\bibitem[Str96]{Str96} Stein Arild Strømme, \emph{Elementary introduction to representable functors and Hilbert schemes}, Parameter spaces (Warsaw, 1994), Banach Center Publ., vol. 36, Polish Acad. Sci., Warsaw, 1996, 第 179--198 页.
\bibitem[Str03]{Str03} Ross Street, \emph{Categorical and combinatorial aspects of descent theory}, arXiv: math.CT/0303175, 2003.
\bibitem[Tam04]{Tam04} Akio Tamagawa, \emph{Finiteness of isomorphism classes of curves in positive characteristic with prescribed fundamental groups}, J. Algebraic Geom. \textbf{13} (2004), no. 4, 675--724.
\bibitem[Tho87]{Tho87} Robert W. Thomason, \emph{Algebraic $K$-theory of group scheme actions}, Algebraic topology and algebraic $K$-theory (Princeton, N.J., 1983), Ann. of Math. Stud., vol. 113, Princeton Univ. Press, Princeton, NJ, 1987, 第 539--563 页.
\bibitem[Vid01]{Vid01} Isabelle Vidal, \emph{Contributions à la cohomologie étale des schémas et des log schémas}, 博士学位论文, Paris XI Orsay, 2001.
\bibitem[Vis97]{Vis97} Angelo Vistoli, \emph{The deformation theory of local complete intersections}, arXiv: alg-geom/9703008.
\bibitem[VW94]{VW94} Cumrun Vafa and Edward Witten, \emph{A strong coupling test of S-duality}, Nuclear Phys. B \textbf{431} (1994), no. 1--2, 3--77.
\bibitem[Wat79]{Wat79} William C. Waterhouse, \emph{Introduction to affine group schemes}, Graduate Texts in Mathematics, vol. 66, Springer-Verlag, New York, 1979.
\bibitem[Wei54]{Wei54} André Weil, \emph{Sur les critères d'équivalence en géométrie algébrique}, Math. Ann. \textbf{128} (1954), 95--127.
\bibitem[Wei79]{Wei79} André Weil, \emph{Commentaire}, Collected papers, vol. Vol. I, 第 518--574 页; Vol. II, 第 526--553 页; Vol. III, 第 443--465 页., Springer-Verlag, 1979.
\bibitem[Zar58]{Zar58} Oscar Zariski, \emph{Introduction to the problem of minimal models in the theory of algebraic surfaces}, Publications of the Mathematical Society of Japan, no. 4, The Mathematical Society of Japan, Tokyo, 1958.
\bibitem[Zar71]{Zar71} \rule{2em}{.3pt}, \emph{Algebraic surfaces}, 增补版, Springer-Verlag, New York, 1971, 附 S. S. Abhyankar, J. Lipman 和 D. Mumford 所作的附录, Ergebnisse der Mathematik und ihrer Grenzgebiete, Band 61.
\bibitem[ZS60]{ZS60} Oscar Zariski and Pierre Samuel, \emph{Commutative algebra. Vol. II}, The University Series in Higher Mathematics, D. Van Nostrand Co., Inc., Princeton, N. J.-Toronto-London-New York, 1960.
% PDF344 / printed332 is blank in the source scan.
\end{thebibliography}
